% Owner: E. M4 (H5), input from methods.tex.
\textbf{M4} cuts the event sequence into consecutive windows of width $\Delta t$ and
aggregates each into a weighted undirected snapshot (weight = number of events). Windows
with very few events are dropped (fewer than 100 for PS, which removes nights; fewer
than 50 for CM), and only nodes active in a window belong to its snapshot. The main
settings are $\Delta t=15$\,min for PS and 7\,days for CM. Every snapshot and the
static aggregate are partitioned with Leiden \citep{traag2019louvain} maximising
modularity with resolution $\gamma$ \citep{reichardt2006statistical}. A multislice
partition \citep{mucha2010community} with interslice coupling $\omega=1$ gives labels
that are consistent across windows; we use it only for tracking, because the coupling
itself pulls consecutive partitions together.

Partitions are compared with NMI \citep{danon2005comparing}, always on the nodes the two
partitions share. For each window we record the NMI with the previous window (only for
back-to-back windows, not across a night), with the static partition, and for PS with
the classes (teachers excluded). As a noise floor, every window is partitioned a second
time with another random seed. All measures are recomputed on 20 P1 samples: P1 keeps
the static graph, and hence the static partition, and the number of events per window,
so H5 predicts that real snapshots agree less with each other and with the static
partition than P1 snapshots do. We repeat the analysis for $\gamma\in\{0.5,1,2\}$ and
$\Delta t\in\{5,15,60\}$\,min (PS) or $\{1,7,30\}$\,days (CM).
